\documentclass[10pt]{article}
\usepackage[margin=20mm]{geometry}
\usepackage{amsmath,amssymb,amsthm}
\newtheorem{lemma}{Lemma}
\newtheorem{proposition}[lemma]{Proposition}
\newcommand{\diag}{\operatorname{diag}}
\newcommand{\off}{\operatorname{off}}
\begin{document}
\begin{center}\large An exact finite Gram cone for quadratic half-line selectors\end{center}

Fix the published canonical source
$c=(d,\lambda,n,m,C,A,D,\alpha,\beta)\ge0$. Let $P(v),N(v),E(v),K(v)$
be the four symmetric $2\times2$ matrices of the canonical-completion note.
The cone considered here consists precisely of sources for which some
$J(v)=J_0+vJ_1+v^2J_2$, with unrestricted real $2\times2$ coefficients,
makes
\[
 G_J(v)=\begin{pmatrix}P+K&J\\J^T&N+K\end{pmatrix},\qquad
 H_J(v)=\begin{pmatrix}P+E&K-E-J\\K-E-J^T&N+E\end{pmatrix}
 \succeq0\quad(v\ge0).                                      \tag{1}
\]
Equality of this cone with all valid canonical sources is not assumed.

\begin{lemma}[Quadratic matrix polynomials on a half-line]
For real symmetric $4\times4$ matrices $M_0,M_1,M_2$, the polynomial
$M(v)=M_0+vM_1+v^2M_2$ is positive semidefinite for all $v\ge0$ if and
only if there are an $8\times8$ Gram matrix $S\succeq0$ and a
$4\times4$ Gram matrix $T\succeq0$ with
\[
 S=\begin{pmatrix}M_0&B\\B^T&M_2\end{pmatrix},\qquad
 T=M_1-B-B^T.                                               \tag{2}
\]
Here $B$ is an unrestricted real $4\times4$ matrix.
\end{lemma}
\begin{proof}
The reverse implication is the identity
$M(v)=[I\ vI]S[I\ vI]^T+vT$.
Conversely $M(t^2)$ is a real matrix polynomial of degree at most four
which is positive semidefinite on the real line. The real-line matrix
Fej\'er--Riesz theorem (Zalar, \emph{Matrix Fej\'er--Riesz theorem with gaps},
Theorem~1.2) gives a polynomial matrix factor of degree at most
two; taking real and imaginary parts gives real matrix squares.
Averaging their values at $t$ and $-t$ yields
\[
 M(v)=(R_0+vR_2)^T(R_0+vR_2)+vR_1^TR_1.
\]
Set $S=[R_0\ R_2]^T[R_0\ R_2]$ and $T=R_1^TR_1$.
The degree bounds are sharp: $M(v)=vI$ needs the half-line term,
and $M(v)=v^2I$ needs a linear factor. No strict positivity, invertibility,
or nonzero leading coefficient is used. The result also holds over real
closed fields by transfer of this finite first-order assertion.
\end{proof}

Here are the evaluated integer coefficient formulas. Write
$D(a,b)=\diag(a,b)$ and $O(a)=\left(\begin{smallmatrix}0&a\\a&0\end{smallmatrix}\right)$.
In the coefficient order $1,v,v^2$,
\begin{align*}
 (P_0,P_1,P_2)&=(D(0,4m+4\beta+2d),O(-4m),D(4m,4\beta+2d)),\\
 (N_0,N_1,N_2)&=(D(4C+4D+2d,0),O(4D+2d),D(4C,4D+2d)),\\
 (E_0,E_1,E_2)&=(D(4A,4\alpha),O(-4\alpha-2d),D(4A+4\alpha+2d,-2d)),\\
 (K_0,K_1,K_2)&=(D(4\lambda,4n),O(4\lambda),D(0,4\lambda+4n)).
\end{align*}
For $j=0,1,2$ define
\[
 G_j=\begin{pmatrix}P_j+K_j&J_j\\J_j^T&N_j+K_j\end{pmatrix},\qquad
 H_j=\begin{pmatrix}P_j+E_j&K_j-E_j-J_j\\
                         K_j-E_j-J_j^T&N_j+E_j\end{pmatrix}.    \tag{3}
\]
These are affine linear with integer coefficients in the source and selector
entries. In addition to the twelve selector entries, introduce unrestricted
$4\times4$ matrices $B_G,B_H$. The complete coefficient elimination is
\[
 \boxed{\begin{aligned}
 S_G&=\begin{pmatrix}G_0&B_G\\B_G^T&G_2\end{pmatrix}\succeq0,
 &S_H&=\begin{pmatrix}H_0&B_H\\B_H^T&H_2\end{pmatrix}\succeq0,\\
 T_G&=G_1-B_G-B_G^T\succeq0,
 &T_H&=H_1-B_H-B_H^T\succeq0.
 \end{aligned}}                                             \tag{4}
\]
\begin{proposition}
Projection of (4) onto $c\ge0$ is exactly the cone in (1), including every
zero or degenerate source face. It has four Gram blocks of orders
$8,8,4,4$, total order $24$, and $12+16+16=44$ auxiliary scalars, with no
remaining coefficient equations. Including the nine scalar blocks $[c_i]$
gives an ambient affine semidefinite representation of total order $33$.
\end{proposition}
\begin{proof}
If (4) holds, (2) reconstructs each matrix in (1) for every $v\ge0$.
If (1) holds, apply the lemma independently to its two quadratic matrices;
their off-diagonal Gram blocks give $B_G,B_H$ in (4). Thus no selector has
been discarded and no extra source condition introduced. All operations
are coefficient identities without division, so zeros and singular Gram
matrices remain literal cases. In particular the all-zero data are feasible.
\end{proof}

For comparison, retaining all Gram entries would give $72+20+12=104$
unknowns and $60$ scalar coefficient equations. Formula (4) solves these
equations explicitly, leaving $44$ unknowns; it makes no symmetry assumption
on either $B$ or any $J_j$.

\begin{lemma}[Bounded auxiliaries and closed projection]
Put $s=\sum_i c_i$. Every feasible auxiliary vector satisfies
$\|y\|_\infty\le16s$, and the projected source cone is closed. More precisely,
\begin{align*}
 g:=\operatorname{tr}S_G&=8d+16\lambda+16n+8m+8C+8D+8\beta,\\
 h:=\operatorname{tr}S_H&=8d+8m+8C+8D+8\beta+16A+16\alpha,
 \qquad \operatorname{tr}T_G\le g,\quad\operatorname{tr}T_H\le h.
\end{align*}
\end{lemma}
\begin{proof}
All diagonals of $S_G,S_H$ depend only on $c$, and $g,h\le16s$.
Their PSD principal minors bound every $B_G,B_H,J_0,J_2$ entry by
$g/2$ or $h/2$. Put $a_i=(S_G)_{ii}+(S_G)_{i+4,i+4}$, $0\le i<4$.
Since $(G_1)_{ii}=0$, positivity gives
$0\le(T_G)_{ii}=-2(B_G)_{ii}\le a_i$; summing gives the trace bound,
and likewise for $H$. For each selector entry $(J_1)_{ij}$, its position
in $G_1$ has distinct indices $r=i$, $t=j+2$. The PSD entry bounds give
$|(J_1)_{ij}|\le |(T_G)_{rt}|+|(B_G)_{rt}|+|(B_G)_{tr}|
\le a_r+a_t\le g$.
This bounds all 44 auxiliaries. Along a convergent sequence of feasible
sources the auxiliaries therefore have a convergent subsequence, whose
limit remains feasible by closedness of the PSD cones. With arbitrary
signed source inputs the same proof uses $16\sum_i|c_i|$, so the order-24
projection is closed as well.
Moreover fixed-source feasibility cannot be weakly infeasible. Otherwise a
sequence approaching the PSD cones would either have bounded auxiliaries,
giving a feasible limit, or an unbounded subsequence. Dividing the latter by
its auxiliary norm and taking a subsequence would give feasible homogeneous
data with source zero and a nonzero auxiliary vector, contradicting the bound.
\end{proof}

For each feasible source, put $T_2=\diag(1,-1)$. The published matrices obey
$M(-v)=T_2M(v)T_2$. Define the selector on $v<0$ by
$J_-(v)=T_2J_+(-v)T_2$. Congruence by $\diag(T_2,T_2)$ carries both
inequalities from $-v>0$ to $v<0$. Thus (4) implies canonical validity by
the pointwise completion theorem; finite axes and projective endpoints
follow from the full homogeneous parent. This argument asserts only this
inclusion. It does not prove that every valid source admits a quadratic
half-line selector. The reflected selector is generally piecewise
polynomial, so the global polynomial-selector counterexample does not
contradict this construction.

The two accompanying COO files evaluate (3)--(4) exactly. They use all
53 ordered parameters (nine source and 44 auxiliary entries), full symmetric
COO indexing, and respectively omit or include the external source-sign
blocks. The standard-library checker reconstructs both original matrix
quadratics for every parameter basis vector at three rational testing
values; hence it checks every coefficient. It performs no solver call,
numerical search, or universal-adequacy test.
\end{document}
